
\documentclass{article}
\usepackage{colortbl}
\usepackage{makecell}
\usepackage{multirow}
\usepackage{supertabular}

\begin{document}

\newcounter{utterance}

\centering \large Interaction Transcript for game `cladder', experiment `full\_v1.5\_default', episode 9197 with reflexion{-}llama{-}3.1{-}8b{-}instruct.
\vspace{24pt}

{ \footnotesize  \setcounter{utterance}{1}
\setlength{\tabcolsep}{0pt}
\begin{supertabular}{c@{$\;$}|p{.15\linewidth}@{}p{.15\linewidth}p{.15\linewidth}p{.15\linewidth}p{.15\linewidth}p{.15\linewidth}}
    \# & \multicolumn{2}{c}{Player} && \multicolumn{2}{c}{Game Master} \\
    \hline

    \theutterance \stepcounter{utterance}  
    & & & \multicolumn{4}{p{0.6\linewidth}}{
        \cellcolor[rgb]{0.9,0.9,0.9}{
            \makecell[{{p{\linewidth}}}]{
                \texttt{\tiny{[P1$\langle$GM]}}
                \texttt{You are an expert in causal inference. The following question is not a typical commonsense query, but rather a meticulously designed question created by a professor specializing in causal inference, intended to assess the students' mastery of the course content.} \\
\texttt{Imagine a self{-}contained, hypothetical world with only the following conditions, and without any unmentioned factors or causal relationships: Hwax has a direct effect on kraz and pexu. Pexu has a direct effect on rukz. Kraz has a direct effect on rukz. The overall probability of pexu is 82\%. The probability of not pexu and rukz is 13\%. The probability of pexu and rukz is 47\%. Is the chance of rukz smaller when observing pexu?} \\
\texttt{Start your answer with 'Yes' or 'No', followed by additional reasoning or evidence to support your explanation.} \\
            }
        }
    }
     \\ \\

    \theutterance \stepcounter{utterance}  
    & \multicolumn{4}{p{0.6\linewidth}}{
        \cellcolor[rgb]{0.561,0.737,0.561}{
            \makecell[{{p{\linewidth}}}]{
                \texttt{\tiny{[P1$\rangle$GM]}}
                \texttt{No.} \\
\\ 
\texttt{To determine whether the chance of rukz is smaller when observing pexu, we need to assess the causal relationships and conditional probabilities in this hypothetical world.} \\
\\ 
\texttt{First, let's establish the causal relationships:} \\
\\ 
\texttt{1. Hwax → kraz} \\
\texttt{2. Hwax → pexu} \\
\texttt{3. pexu → rukz} \\
\texttt{4. kraz → rukz} \\
\\ 
\texttt{We are given the following probabilities:} \\
\\ 
\texttt{{-} P(pexu) = 82\%} \\
\texttt{{-} P(¬pexu ∧ rukz) = 13\%} \\
\texttt{{-} P(pexu ∧ rukz) = 47\%} \\
\\ 
\texttt{We can calculate the probability of rukz as follows:} \\
\\ 
\texttt{P(rukz) = P(¬pexu ∧ rukz) + P(pexu ∧ rukz)} \\
\texttt{= 13\% + 47\%} \\
\texttt{= 60\%} \\
\\ 
\texttt{Now, we need to determine whether the presence of pexu affects the probability of rukz. To do this, we can use the law of total probability:} \\
\\ 
\texttt{P(rukz | pexu) = P(rukz ∧ pexu) / P(pexu)} \\
\texttt{= 47\% / 82\%} \\
\texttt{= 57.32\%} \\
\\ 
\texttt{P(rukz | ¬pexu) = P(rukz ∧ ¬pexu) / P(¬pexu)} \\
\texttt{= P(¬pexu ∧ rukz) / (1 {-} P(pexu))} \\
\texttt{= 13\% / (1 {-} 82\%)} \\
\texttt{= 13\% / 18\%} \\
\texttt{= 72.22\%} \\
\\ 
\texttt{Comparing the conditional probabilities, we see that P(rukz | pexu) = 57.32\% is less than P(rukz | ¬pexu) = 72.22\%. However, this does not directly answer the question. We need to consider the effect of pexu on rukz, taking into account the direct and indirect effects of hwax on rukz.} \\
\\ 
\texttt{Since pexu has a direct effect on rukz, and kraz also has a direct effect on rukz, we can use the formula for the total effect of a common cause:} \\
\\ 
\texttt{P(rukz | pexu) = P(rukz | pexu) + P(rukz | kraz) {-} P(rukz | pexu, kraz)} \\
\\ 
\texttt{However, we do not have the necessary information to calculate this expression. Nevertheless, we can use the given probabilities to assess the relationship between pexu and rukz.} \\
\\ 
\texttt{The presence of pexu increases the probability of rukz, but the effect is not as strong as the effect of kraz on rukz. This suggests that the relationship between pexu and rukz is not a direct causal relationship, but rather an indirect effect mediated by kraz.} \\
\\ 
\texttt{In conclusion, the presence of pexu does not decrease the probability of rukz. Instead, it increases the probability of rukz, but} \\
            }
        }
    }
    & & \\ \\

    \theutterance \stepcounter{utterance}  
    & & & \multicolumn{2}{p{0.3\linewidth}}{
        \cellcolor[rgb]{0.9,0.9,0.9}{
            \makecell[{{p{\linewidth}}}]{
                \texttt{\tiny{[GM$|$GM]}}
                \texttt{No} \\
            }
        }
    }
    & & \\ \\

    \theutterance \stepcounter{utterance}  
    & & & \multicolumn{2}{p{0.3\linewidth}}{
        \cellcolor[rgb]{0.9,0.9,0.9}{
            \makecell[{{p{\linewidth}}}]{
                \texttt{\tiny{[GM$|$GM]}}
                \texttt{yes} \\
            }
        }
    }
    & & \\ \\

    \theutterance \stepcounter{utterance}  
    & & & \multicolumn{2}{p{0.3\linewidth}}{
        \cellcolor[rgb]{0.9,0.9,0.9}{
            \makecell[{{p{\linewidth}}}]{
                \texttt{\tiny{[GM$|$GM]}}
                \texttt{game\_result = LOSE} \\
            }
        }
    }
    & & \\ \\

    \theutterance \stepcounter{utterance}  
    & & \multicolumn{4}{p{0.6\linewidth}}{
        \cellcolor[rgb]{0.95,0.95,0.95}{
            \makecell[{{p{\linewidth}}}]{
                \texttt{\tiny{[P1$\leftrightarrow$P]}}
                \texttt{Post{-}episode Reflexion (not sent to GM):} \\
\texttt{New plan: } \\
\\ 
\texttt{In my previous attempt, I correctly identified the causal relationships and conditional probabilities in the hypothetical world. However, I failed to properly assess the effect of pexu on rukz, taking into account the direct and indirect effects of hwax on rukz. } \\
\\ 
\texttt{To improve my strategy, I will focus on correctly applying the formula for the total effect of a common cause, which is:} \\
\\ 
\texttt{P(rukz | pexu) = P(rukz | pexu) + P(rukz | kraz) {-} P(rukz | pexu, kraz)} \\
\\ 
\texttt{I will also pay closer attention to the given probabilities and ensure that I use them correctly to assess the relationship between pexu and rukz. Specifically, I will re{-}examine the probabilities of pexu and rukz, and consider how the presence of pexu affects the probability of rukz, taking into account the direct and indirect effects of hwax on rukz.} \\
\\ 
\texttt{To achieve this, I will break down the problem into smaller, more manageable steps, and carefully apply the relevant formulas and concepts from causal inference. I will also double{-}check my calculations and reasoning to ensure that I arrive at a correct and well{-}supported conclusion.} \\
            }
        }
    }
    & \\ \\

\end{supertabular}
}

\end{document}
